
\chapter{Implementación del pipeline CI/CD}

\section{Plataforma utilizada}

\campo{
Indicar la plataforma utilizada (GitLab, GitHub Actions,
Azure DevOps u otra) y justificar su elección.
}

\section{Arquitectura del pipeline}

\campo{
Describir el flujo de integración y entrega continua,
identificando sus etapas y dependencias.
}

\evidencefigure
{figures/pipeline.pdf}
{Diagrama general del pipeline CI/CD.}
{fig:pipeline}

\section{Etapas del pipeline}

\begin{table}[H]
\centering
\caption{Etapas y trabajos del pipeline.}
\label{tab:stages}

\begin{tabularx}{\textwidth}{l l X}
\toprule
\textbf{Etapa} &
\textbf{Job} &
\textbf{Descripción} \\
\midrule
\campo{Stage 1} & \campo{Job} & \campo{Descripción} \\
\campo{Stage 2} & \campo{Job} & \campo{Descripción} \\
\campo{Stage 3} & \campo{Job} & \campo{Descripción} \\
\campo{Stage 4} & \campo{Job} & \campo{Descripción} \\
\bottomrule
\end{tabularx}
\end{table}

\section{Configuración del pipeline}

\campo{
Presentar el archivo de configuración y explicar los elementos
más relevantes de su implementación.
}

\subsection{Fragmentos de configuración}

\begin{lstlisting}[
    language=yaml,
    caption={Ejemplo de estructura del pipeline.},
    label={lst:pipeline}
]
# Sustituir por fragmentos del pipeline real

stages:
  - build
  - test
  - security
  - deploy

# Agregar aqui los jobs implementados
\end{lstlisting}

\section{Evidencia de ejecución}

\campo{
Incluir capturas del pipeline ejecutado, estados de los jobs
y generación de artefactos.
}

\evidencefigure
{evidence/pipeline/ejecucion.png}
{Ejecución del pipeline CI/CD.}
{fig:pipeline-ejecucion}

\section{Repositorio y estructura base}

\begin{itemize}
    \item \textbf{Repositorio del proyecto:}

    \externallink{https://github.com/USUARIO/REPOSITORIO}
    {Enlace al repositorio}

    \item \textbf{Estructura base del pipeline:}

    \externallink{https://gitlab.com/-/snippets/5999866}
    {Plantilla de jobs DevSecOps}

    \item \textbf{Archivo de configuración:}

    \campo{Ruta del archivo .gitlab-ci.yml o equivalente.}
\end{itemize}
